\begin{ouverture}[milieux]{14}{La droite des milieux}{Partie 4 · Géométrie et espace}
{Milieux · Parallèles · Moitié · Démonstrations}
\objectif{Connaître les trois théorèmes relatifs à la droite des milieux.}
\objectif{Choisir le théorème adapté pour prouver un parallélisme, un milieu ou calculer une longueur.}
\objectif{Rédiger une démonstration en plusieurs étapes.}
\prerequis{Demi-tour et parallélogrammes (Cinquième, chapitre 13).}
\end{ouverture}
\pagestyle{otMain}\essentiel\fontsize{11}{15}\selectfont
\phantomsection\label{cours-14-1}
\essSection{1}{Trois théorèmes}
\begin{theoreme}[1 · Droite des milieux]Dans un triangle, la droite qui passe par les milieux de deux côtés est parallèle au troisième côté.\end{theoreme}
\begin{theoreme}[2 · Segment des milieux]Dans un triangle, le segment qui joint les milieux de deux côtés mesure la moitié du troisième côté.\end{theoreme}
\begin{theoreme}[3 · Parallèle passant par un milieu]Dans un triangle, la droite qui passe par le milieu d’un côté et qui est parallèle à un deuxième côté coupe le troisième côté en son milieu.\end{theoreme}
\label{fig-cours-14-1-cours-milieux}\figureGeo{cours-milieux}

\begin{exemple}Dans le triangle ABC, I est le milieu de $[AB]$ et J celui de $[AC]$, avec $BC=9$ cm. D’après les théorèmes 1 et 2, $(IJ)\parallel(BC)$ et $IJ=\frac92=4{,}5$ cm.\end{exemple}
\clearpage
\phantomsection\label{cours-14-2}
\essSection{2}{Démontrer et choisir}
\begin{demonstration}[Idée de la preuve des théorèmes 1 et 2]Le demi-tour de centre J transforme I en un point K. Les diagonales $[AC]$ et $[IK]$ de AICK ont le même milieu J : c’est un parallélogramme, donc $(CK)\parallel(IB)$ et $CK=AI=IB$. Alors IBCK, qui a deux côtés opposés parallèles et de même longueur, est un parallélogramme : $(IJ)\parallel(BC)$ et $IJ=\frac12IK=\frac12BC$. La preuve détaillée, avec le théorème 3, figure dans « Comprendre les preuves ».\end{demonstration}
\begin{savoirfaire}[Choisir le théorème]Pour prouver un parallélisme à partir de deux milieux : théorème 1. Pour calculer une longueur à partir de deux milieux : théorème 2. Pour prouver qu’un point est un milieu à partir d’un milieu et d’un parallélisme : théorème 3.\end{savoirfaire}
\begin{savoirfaire}[Rédiger]« Dans le triangle ABC, I est le milieu de $[AB]$ et J est le milieu de $[AC]$. D’après le théorème de la droite des milieux, $(IJ)$ est parallèle à $(BC)$. » On cite le triangle, les données, le théorème, puis la conclusion.\end{savoirfaire}
\begin{attention}Les milieux doivent être donnés par l’énoncé ou démontrés, jamais lus sur le dessin. Le segment des milieux mesure la moitié du troisième côté : celui qui ne contient aucun des deux milieux.\end{attention}

\clearpage\rubrique{Découverte}
\begin{blocLivre}{ex-14-1}\begin{exercice}[1]{Une longueur}
Dans le triangle ABC, I est le milieu de $[AB]$, J celui de $[AC]$ et $BC=7{,}4$ cm. Calculer IJ.
\label{fig-ex-14-1}\figureGeo{mil-01}
\end{exercice}\end{blocLivre}
\begin{blocLivre}{ex-14-2}\begin{exercice}[1]{Un parallélisme}
Dans le triangle EFG, M est le milieu de $[EF]$ et N celui de $[EG]$. Démontrer que $(MN)$ et $(FG)$ sont parallèles.
\label{fig-ex-14-2}\figureGeo{mil-02}
\end{exercice}\end{blocLivre}
\begin{blocLivre}{ex-14-3}\begin{exercice}[1]{Retrouver un côté}
Dans le triangle ABC, I et J sont les milieux respectifs de $[AB]$ et $[AC]$, et $IJ=5{,}2$ cm. Calculer BC.
\label{fig-ex-14-3}\figureGeo{mil-03}
\end{exercice}\end{blocLivre}
\begin{blocLivre}{ex-14-4}\begin{exercice}[1]{Un milieu}
Dans le triangle ABC, I est le milieu de $[AB]$. La parallèle à $(BC)$ passant par I coupe $[AC]$ en J. Sachant que $AC=9$ cm, calculer AJ.
\label{fig-ex-14-4}\figureGeo{mil-04}
\end{exercice}\end{blocLivre}
\begin{blocLivre}{ex-14-5}\begin{exercice}[1]{Choisir le théorème}
Dans chaque figure, les codages sont des données. Indiquer le théorème qui permet de conclure, puis conclure.\begin{questions}
\item Figure 1 : les droites $(MN)$ et $(FG)$ sont-elles parallèles ?
\item Figure 2 : où se trouve le point N sur $[EG]$ ?
\item Figure 3 : sachant que $FG=8$ cm, que vaut MN ?
\end{questions}
\label{fig-ex-14-5}\figureGeo{mil-05}
\end{exercice}\end{blocLivre}

\rubrique{Entraînement}
\begin{blocLivre}{ex-14-6}\begin{exercice}[2]{Le triangle des milieux}
Le triangle ABC a pour côtés $AB=6$ cm, $AC=8$ cm et $BC=10$ cm. I, J et K sont les milieux de $[AB]$, $[AC]$ et $[BC]$. Calculer le périmètre du triangle IJK et le comparer à celui de ABC.
\label{fig-ex-14-6}\figureGeo{mil-06}
\end{exercice}\end{blocLivre}
\begin{blocLivre}{ex-14-7}\begin{exercice}[2]{Un milieu, puis une longueur}
Dans le triangle ABC, I est le milieu de $[AB]$ et la parallèle à $(BC)$ passant par I coupe $[AC]$ en J. On sait que $BC=9$ cm. Calculer IJ en justifiant chaque étape.
\label{fig-ex-14-7}\figureGeo{mil-07}
\end{exercice}\end{blocLivre}
\begin{blocLivre}{ex-14-8}\begin{exercice}[2]{Une rédaction incomplète}
Un élève écrit : « I est le milieu de $[AB]$ et J est un point de $[AC]$, donc $(IJ)$ est parallèle à $(BC)$. » Pourquoi ce raisonnement est-il incorrect ? Que faudrait-il savoir ?
\label{fig-ex-14-8}\figureGeo{mil-08}
\end{exercice}\end{blocLivre}
\begin{blocLivre}{ex-14-9}\begin{exercice}[2]{Dans un rectangle}
ABCD est un rectangle de diagonale $AC=10$ cm. I est le milieu de $[AB]$ et J celui de $[BC]$. Démontrer que $(IJ)$ est parallèle à $(AC)$ et calculer IJ.
\label{fig-ex-14-9}\figureGeo{mil-09}
\end{exercice}\end{blocLivre}
\begin{blocLivre}{ex-14-10}\begin{exercice}[2]{Un parallélogramme de milieux}
Dans le triangle ABC, I, J et K sont les milieux de $[AB]$, $[AC]$ et $[BC]$, avec $AB=7$ cm et $AC=5$ cm. Démontrer que AIKJ est un parallélogramme, puis calculer son périmètre.
\label{fig-ex-14-10}\figureGeo{mil-10}
\end{exercice}\end{blocLivre}

\rubrique{Problèmes}
\begin{blocLivre}{ex-14-11}\begin{exercice}[3]{Des milieux de milieux}
Dans le triangle ABC, $BC=12$ cm. I et J sont les milieux de $[AB]$ et $[AC]$, puis M et N sont les milieux de $[AI]$ et $[AJ]$. Calculer MN et justifier que $(MN)\parallel(BC)$.
\label{fig-ex-14-11}\figureGeo{mil-11}
\end{exercice}\end{blocLivre}
\begin{blocLivre}{ex-14-12}\begin{exercice}[3]{Dans un trapèze}
ABCD est un trapèze de bases $[AB]$ et $[CD]$, avec $AB=6$ cm et $CD=10$ cm. I est le milieu de $[AD]$. La parallèle à $(AB)$ passant par I coupe $[BD]$ en J et $[BC]$ en K. Démontrer que J et K sont des milieux, puis calculer IJ, JK et IK.
\label{fig-ex-14-12}\figureGeo{mil-12}
\end{exercice}\end{blocLivre}
\begin{blocLivre}{ex-14-13}\begin{exercice}[3]{Le théorème de Varignon}
ABCD est un quadrilatère non croisé quelconque. I, J, K et L sont les milieux de $[AB]$, $[BC]$, $[CD]$ et $[DA]$. Démontrer que IJKL est un parallélogramme. Calculer son périmètre si $AC=8$ cm et $BD=6$ cm.
\label{fig-ex-14-13}\figureGeo{mil-13}
\end{exercice}\end{blocLivre}
\begin{blocLivre}{ex-14-14}\begin{exercice}[3]{Le point de concours des médianes}
Dans le triangle ABC, I est le milieu de $[BC]$ et J celui de $[AC]$. Les médianes $(AI)$ et $(BJ)$ se coupent en G. K et L sont les milieux de $[AG]$ et $[BG]$. Démontrer que KLIJ est un parallélogramme. En déduire AG si $AI=9$ cm.
\label{fig-ex-14-14}\figureGeo{mil-14}
\end{exercice}\end{blocLivre}
\begin{blocLivre}{ex-14-15}\begin{exercice}[3]{Le milieu de l’hypoténuse}
ABC est rectangle en A, I est le milieu de $[BC]$ et J celui de $[AB]$. Démontrer que $(IJ)$ est la médiatrice de $[AB]$. En déduire que $IA=IB=IC$ et calculer IA si $BC=10$ cm.
\label{fig-ex-14-15}\figureGeo{mil-15}
\end{exercice}\end{blocLivre}

\clearpage
